\documentclass[11pt,a4paper]{article}
\usepackage[margin=27mm]{geometry}
\usepackage{amsmath,amssymb}
\usepackage[T1]{fontenc}
\setlength{\emergencystretch}{2em}
\title{Summable inexact-inclusion errors do not ensure convergence\\
\large Disproof of conjecture 00000008835}
\author{ziangni-sys}
\date{4 October 2026}
\begin{document}
\maketitle
The conjecture asserts that summable tolerance sequences give
strong convergence for approximate inclusions. Both language
versions leave this assertion unqualified.
An exact proximal-point iteration is already a counterexample:
the tolerances are identically zero, but the operator has no zero.
This refutes the convergence conjunct; the separate, unspecified
product-rate assertion is not needed.

\textbf{The operator.}
Use the real Hilbert space $\mathbb R$ with its usual inner product,
and the set-valued operator
\[
 A(x)=\{1\}\qquad(x\in\mathbb R).
\]
Its graph is $G=\{(x,1):x\in\mathbb R\}$.
It is monotone, since for any two graph points
\[
 (x-y)(1-1)=0\geq0.
\]
It is also maximal monotone. Indeed, let a monotone graph $H$
contain $G$, and take $(u,v)\in H$. Monotonicity with
$(u-1,1)$ and $(u+1,1)$ gives respectively
\[
 v-1\geq0,\qquad -(v-1)\geq0.
\]
Thus $v=1$, so $(u,v)\in G$. No proper monotone enlargement exists.

The single-valued representative $x\mapsto1$ is globally
Lipschitz: its actual Lipschitz modulus is zero, and in particular
it is $1$-Lipschitz. Nevertheless
\[
 \operatorname{zer} A=\varnothing.
\]
No degeneracy of the proximal parameter or failure of maximality
is involved.

\textbf{Actual resolvent.}
With unit proximal parameter, the defining resolvent equation is
\[
 x\in y+A(y).
\]
Since $A(y)=\{1\}$, this is equivalent to $x=y+1$.
It therefore has the unique solution
\[
 J_A(x)=(I+A)^{-1}(x)=x-1
\]
for every $x\in\mathbb R$.

\newpage
\textbf{Genuine approximate-inclusion steps.}
The standard additive-error proximal inclusion with tolerances is
\[
 u_n-u_{n+1}-e_n\in A(u_{n+1}),\qquad
 |e_n|\leq\varepsilon_n .
\]
Take an arbitrary initial point $u_0=x$ and iterate the actual
resolvent:
\[
 u_{n+1}=J_A(u_n),\qquad e_n=0,\qquad\varepsilon_n=0.
\]
An induction gives the complete orbit
\[
 u_n=x-n.
\]
Consequently $u_n-u_{n+1}-e_n=1\in A(u_{n+1})$ and
$|e_n|=0\leq\varepsilon_n$ for every $n$.
The tolerances are nonnegative and summable, with
\[
 \sum_{n=0}^{\infty}\varepsilon_n=0.
\]
This is an exact method within the approximate-inclusion
framework, rather than an arbitrary sequence satisfying a
postulated convergence certificate.

\textbf{Failure of strong convergence.}
The orbit has no finite real limit for any initial point.
For a direct argument, if $u_n\to a$, then also $u_{n+1}\to a$,
so $u_{n+1}-u_n\to0$. But the actual difference is constantly
$-1$. Uniqueness of limits would give $-1=0$, a contradiction.
Thus summability, even zero total error, does not ensure strong
convergence.

\textbf{Scope and the missing hypothesis.}
The counterexample does not contradict a proximal-point theorem
that assumes a nonempty zero set. Solvability is precisely absent
from the source's unconditional convergence assertion.
The operator here is everywhere defined, maximal monotone and
Lipschitz, and every proximal subproblem has a unique solution.
The only essential obstruction is the absence of a target zero.
No assertion is made about the source's vague product of a
summable rate and a Lipschitz constant.

\textbf{Formal correspondence and reproduction.}
Lean defines the actual set-valued graph, the Hilbert inner-product
monotonicity and maximality conditions, and the defining
resolvent relation. It proves the unique formula $J_A(x)=x-1$,
all iterates, the complete inexact-inclusion and norm-error
conditions, nonnegative summable tolerances with sum zero,
and nonconvergence in the real norm topology for every start.

Run \texttt{lake build} in \texttt{lean/} using Lean 4.19.0.
Public Mathlib revision:
\begin{center}\small
\texttt{c44e0c8ee63ca166450922a373c7409c5d26b00b}
\end{center}
The project prints the relevant axiom audits.
\end{document}
